\documentclass{resume}

\usepackage[left=0.75in,top=0.55in,right=0.75in,bottom=0.55in]{geometry}
\usepackage{xcolor}
\usepackage{hyperref}
\usepackage{changepage}
\usepackage{fontawesome}
\usepackage{tabularx}
\usepackage{enumitem}
\usepackage{microtype}

\definecolor{navyblue}{RGB}{0,54,123}
\definecolor{lightgray}{RGB}{245,245,245}

\hypersetup{
  colorlinks=true,
  urlcolor=navyblue,
  linkcolor=navyblue
}

\newcommand{\link}[2]{{\href{#1}{#2}}}
\newcommand{\entry}[2]{#1 & #2 \tabularnewline}

\newcommand{\tableEnv}[2]{%
  \begin{rSection}{#1}
    \begin{adjustwidth}{0.0in}{0.1in}
      \begin{tabularx}{\linewidth}{@{} >{\bfseries}l @{\hspace{5ex}} X @{}}
        #2
      \end{tabularx}
    \end{adjustwidth}
  \end{rSection}
}

\name{\color{navyblue} Kai Ma}

\begin{document}

\printPersonalInfo{
  \personalInfo{%
    \tag{\faEnvelopeO}\info{\href{mailto:makai2024@shanghaitech.edu.cn}{makai2024@shanghaitech.edu.cn}}
    \hspace{1em}
    \tag{\faGraduationCap}\info{\href{https://scholar.google.com/citations?user=hr99haQAAAAJ&hl=en}{Google Scholar}}
    \hspace{1em}
    \tag{\faBarcode}\info{\href{https://orcid.org/0009-0001-4222-9749}{ORCID}}
    \hspace{1em}
    \tag{\faGlobe}\info{\href{https://badbeers.github.io/}{Homepage}}
  }
}

%--- EDUCATION -------------------------------------------------------
\begin{rSection}{Education}

\begin{rSubsectionNoBullet}%
  {\bf ShanghaiTech University}{China}%
  {M.S. in Electronic Science and Technology}{Sep.\ 2024 -- Present}
    \item \textit{Focus: Physical design; Design for manufacturing}
    \item \textit{Supervisor: \href{https://true-genghao.github.io/genghao/}{Dr.\ Hao Geng}}
\end{rSubsectionNoBullet}

\vspace{0.25em}

\begin{rSubsectionNoBullet}%
  {\bf Central China Normal University}{China}%
  {B.S. in Electronic Science and Technology}{Sep.\ 2020 -- Jun.\ 2024}
    \item \textit{Focus: Single-photon detector with low time jitter; Neural network optimization}
\end{rSubsectionNoBullet}

\end{rSection}

%--- PUBLICATIONS ----------------------------------------------------
\begin{rSection}{Publications}
\vspace{0.3em}

\begin{itemize}[label={}, leftmargin=2.5em, itemsep=0.4em]
  \item[{[\textbf{C1}]}] {\bf Kai Ma}, Zhen Wang, Hongquan He, Qi Xu, Tinghuan Chen, Hao Geng,
    ``LMM-IR: Large-Scale Netlist-Aware Multimodal Framework for Static IR-Drop Prediction,''
    \textit{ACM/IEEE Design Automation Conference} ({\bf DAC}), San Francisco, Jun.\ 2025.
  \item[{[\textbf{C2}]}] {\bf Kai Ma}, Tianyi Li, Jiaqi Liu, Jingyi Yu, Hao Geng,
    ``No Pixel, More Efficient: Multimodal Framework for Sub-nm Mask Process Correction,''
    \textit{IEEE/ACM Design, Automation and Test in Europe} ({\bf DATE}), Verona, Italy, Apr.\ 2026.
\end{itemize}

\vspace{0.15em}

\begin{itemize}[label={}, leftmargin=2.5em, itemsep=0.4em]
  \item[{[\textbf{J1}]}] Tianyi Li, {\bf Kai Ma}(co-first author), Yiwen Wu, Jiaqi Liu, Jingyi Yu, Hao Geng,
    ``Shot Splatting: Pioneering Shot-aware MPC for VSB Mask Data Preparation,''
    \textit{IEEE Transactions on Computer-Aided Design of Integrated Circuits and Systems}
    ({\bf TCAD}), doi: 10.1109/TCAD.2026.3667094.
\end{itemize}

\vspace{0.15em}

\begin{itemize}[label={}, leftmargin=2.5em, itemsep=0.4em]
  \item[{[\textbf{C3}]}] Wenxuan Dong, Tianyi Li, {\bf Kai Ma}(co-first author), Zhen Wang, Hanbin Dong, Nan Wang, Su Zheng, Xinyun Zhang, Jingyi Yu, Hao Geng,
    ``UniMRC: A Differentiable Multimodal Framework for Unified Curvilinear-Manhattan Mask Rule Checking,''
    \textit{IEEE/ACM International Conference on Computer-Aided Design} ({\bf ICCAD}), San Jose, Nov.\ 2026.

  \vspace{0.15em}

  \item[{[\textbf{C4}]}] Shenshuo Yao, Liuke Wang, Yuyang Chen, Yinuo Zhu, Zhen Wang, Jiajie Zhang, Tianyi Li, Yiwen Wu, Xinheng Li, {\bf Kai Ma}, Nan Wang, Wenxuan Dong, Xinyun Zhang, Jingyi Yu, Hao Geng,
    ``Ariadne's Thread: A Multimodal Benchmark Weaving Point Solutions into Process Chains for AI-Driven Computational Lithography,''
    \textit{IEEE/ACM International Conference on Computer-Aided Design} ({\bf ICCAD}), San Jose, Nov.\ 2026.
\end{itemize}

\end{rSection}

%--- PATENT ----------------------------------------------------------
\begin{rSection}{Patent}
  Single-photon detector for visible and near-infrared light in the short wavelength
  range with high count rate.
  CN~216668998~U\@. Han Xue et al., \textbf{Kai Ma}, Qinglin Wu.
\end{rSection}

%--- OPEN SOURCE ----------------------------------------------------
\begin{rSection}{Open-Source Repositories}
  \href{https://github.com/TorchOPC/TorchLitho}{\textbf{TorchLitho}} ---
  Contributor to an open-source differentiable computational lithography framework
  (\textbf{226~stars}); implements the first open-source differentiable Abbe / Hopkins
  lithography imaging model and analytical resist development model,
  enabling GPU-accelerated, gradient-based optimisation of Resolution Enhancement
  Techniques (RETs). \textit{(SPIE Advanced Lithography, 2024)}
\end{rSection}

%--- RESEARCH EXPERIENCE ---------------------------------------------
\begin{rSection}{Research Experience}

\begin{rSubsection}%
  {Multimodal EDA: IR-Drop Prediction \& Mask Process Correction}{Sep.\ 2023 -- Present}%
  {Supervisor: \href{https://true-genghao.github.io/genghao/}{Dr.\ Hao Geng}}{ShanghaiTech University, China}
  \item Proposed \textbf{LMM-IR}, a multimodal framework for static IR-drop prediction that encodes
        large-scale SPICE netlists ($>$100K nodes) as 3D point clouds via a novel Large-scale Netlist
        Transformer (LNT), then fuses them with circuit image features through cross-attention;
        achieved the best F1 score and lowest MAE among ICCAD~2023 contest winners. \textit{(DAC 2025)}
  \item Developed a multimodal \textbf{Mask Process Correction (MPC)} framework that represents GDS
        mask layouts as sparse point clouds and conditions correction on e-beam lithography resist
        images via a ControlNet-inspired architecture; achieved sub-nm EPE accuracy comparable
        to commercial tools at $>$200$\times$ speedup on large-scale layouts. \textit{(DATE 2026)}
  \item Proposed \textbf{Shot Splatting}, a unified MPC and shot-fracturing framework; jointly optimises shot geometry, position, size, and dose via
        differentiable Lam\'{e}-curve kernels and EBL simulation, eliminating separate MRC
        post-processing and reducing EPE by $\sim$79--86\% over the no-MPC baseline. \textit{(TCAD)}

  \item Proposed \textbf{OmniShot}, a shot-centric generative framework for Mask-Wafer
        Co-Optimization (MWCO) that unifies OPC, MPC, and shot fracturing into a single
        inference step; developed a CUDA-accelerated soft rasterizer and differentiable EBL
        \& optical lithography simulators for physics-aware end-to-end training; achieved
        $>$40\% lower EPE, $>$40\% fewer shots, and $5\times$ faster runtime over strong
        baselines on complex benchmarks and a full-chip DNN accelerator.
\end{rSubsection}


\begin{rSubsection}%
  {Low Time Jitter Single-Photon Detector}{Sep.\ 2021 -- Mar.\ 2023}%
  {Supervisor: Prof.\ Qinglin Wu}{Central China Normal University, China}
  \item Developed a novel constant fraction timing discrimination method to optimise
        single-photon detection time jitter, achieving a \textbf{50\%} reduction.
  \item Designed an innovative signal processing circuit and conducted end-to-end
        performance validation.
\end{rSubsection}

% \begin{rSubsection}%
%   {Neural Network Optimisation for Filter Design}{Oct.\ 2022 -- May 2023}%
%   {Advisor: Dr.\ Jing Jin}{Central China Normal University, China}
%   \item Applied vector fitting to generate high-quality training and test data,
%         ensuring accuracy across diverse filter configurations.
%   \item Implemented a back-deduction module in NeuroModeler\,Plus via the
%         quasi-Newton algorithm, achieving precise control over transfer-function
%         coefficient ordering.
%   \item Collaborated closely with the research team to fine-tune the optimisation
%         pipeline, delivering efficient and accurate filter designs for multiple
%         applications.
% \end{rSubsection}

\end{rSection}

%--- HONOURS & SCHOLARSHIPS ------------------------------------------
\tableEnv{Awards \& Certifications}{
  \entry{Sep. 2024}{Graduate Academic Scholarship, ShanghaiTech University}
  
  \entry{Jul. 2024}{Digital Implementation \& Sign-off Full-Flow Certification,
    Cadence Design Systems --- delivered by \textit{Miao Liu},
    Sr.\ Group Director.}
    
  \entry{Nov. 2022}{National Encouragement Scholarship, Ministry of Education, China 
  
  \em{— top 5\% of undergraduates nationally}}
  \entry{Nov. 2021}{National Encouragement Scholarship, Ministry of Education, China 
  
  \em{— top 5\% of undergraduates nationally}}
  \entry{Dec. 2021}{Chinese Mathematics Competitions (CMC) --- Second Prize}
}

%--- TECHNICAL SKILLS ------------------------------------------------
\tableEnv{Technical Skills}{
  \entry{Programming}{Python \textperiodcentered\ C/C++ \textperiodcentered\ MATLAB \textperiodcentered\ \LaTeX \textperiodcentered\ Linux}
  \entry{Expertise}{Deep Learning \textperiodcentered\ Virtuoso \textperiodcentered\ Numerical Analysis}
}

%--- CONFERENCE ATTENDANCE -------------------------------------------
\tableEnv{Conference Attendance}{
  \entry{DAC 2025}{ACM/IEEE Design Automation Conference, San Francisco, USA, Jun.\ 2025({\bf Talk \& Poster})}
  \entry{DATE 2026}{IEEE/ACM Design, Automation and Test in Europe, Verona, Italy, Apr.\ 2026({\bf Talk \& Poster})}
  \entry{IWAPS 2025}{International Workshop on Advanced Patterning Solutions, Shenzhen, China, 2025.
    Attended sessions on advanced lithography and mask patterning technologies.}
}

\end{document}